\uuid{7k2A}
\titre{ Réduction d'un endomorphisme de $\mathbb{R}^3$ }

\niveau{}

\module{ Réduction et diagonalisation }


\chapitre{Application linéaire}
\sousChapitre{Diagonalisation}
\theme{}
\auteur{Quentin Liard}
\datecreate{ 2026-03-18}
\organisation{}
\difficulte{}

\datecreate{ 2026-03-05}

\contenu{

\texte{
On considère l’endomorphisme $f$ de $\mathbb{R}^3$ dont la matrice dans la base canonique est

\[
A=
\begin{pmatrix}
2 & 1 & 0\\
1 & 2 & 0\\
0 & 0 & 3
\end{pmatrix}.
\]
}

\begin{enumerate}

\item \question{Calculer le polynôme caractéristique de $A$.}

\item \question{Déterminer les valeurs propres de $A$.}

\item \question{Déterminer une base de chaque sous-espace propre.}

\item \question{Montrer que $A$ est diagonalisable.}

\item \question{Déterminer une matrice diagonale $D$ semblable à $A$ ainsi qu’une matrice de passage $P$.}

\item \question{En déduire l’expression de $A^n$ pour tout $n\in\mathbb{N}$.}

\end{enumerate}

}